%\RequirePackage[l2tabu, orthodox]{nag}
\RequirePackage{currfile}
\documentclass[12pt]{beamer}
\graphicspath{{Imagenes/}{../Imagenes/}}
\usepackage[utf8]{inputenc}
\usepackage[spanish]{babel}
\usepackage{standalone}
\usepackage{color}
\usepackage[binary-units=true]{siunitx}
\usepackage{hyperref}
\hypersetup{
  colorlinks=true,
  linkcolor=blue,          % color of internal links (change box color with linkbordercolor)
  citecolor=green,        % color of links to bibliography
  filecolor=magenta,      % color of file links
  urlcolor=cyan,           % color of external links
  linkbordercolor={0 0 1}
}
\usepackage{xcolor, soul}
\usepackage{etoolbox}
\usepackage{amsmath}
\usepackage{amsthm}
\usepackage{physics}
\usepackage{multicol}
\usepackage{graphicx}
\usepackage{bookmark}
\usepackage{longtable}
\usepackage{graphicx}
\usepackage{tikz}
\usepackage[siunitx, RPvoltages]{circuitikz}
\usetikzlibrary{mindmap}
\usetikzlibrary{arrows, patterns, shapes, decorations.markings, decorations.pathmorphing}
\usetikzlibrary{matrix,positioning}
\tikzstyle{every picture}+=[remember picture,baseline]
\usepackage[autostyle,spanish=mexican]{csquotes}
\usepackage{pifont}
\usepackage[font=footnotesize,textfont=it]{caption}
\usepackage{tabulary}
\usepackage{booktabs}
\usepackage[outdir=./]{epstopdf}
%\usepackage{epstopdf}
\usepackage{media9}
\usepackage{multimedia}
\usepackage{bigints}
%\usepackage{enumitem}
\usepackage[os=win]{menukeys}
\usepackage{pifont}
\usepackage{pbox}
\usepackage{alltt}
\usepackage{verbatim}
\usepackage{colortbl}
\usepackage{tcolorbox}
\usepackage{fancyvrb}
\usepackage[sfdefault]{roboto}  %% Option 'sfdefault' only if the base font of the document is to be sans serif
%\usepackage[T1]{fontenc}
\setcounter{secnumdepth}{3}
\setcounter{tocdepth}{3}
\DeclareGraphicsExtensions{.pdf,.png,.jpg}
\renewcommand {\arraystretch}{1.5}
\definecolor{ao}{rgb}{0.0, 0.5, 0.0}
\definecolor{aquamarine}{rgb}{0.5, 1.0, 0.83}
\definecolor{kellygreen}{rgb}{0.3, 0.73, 0.09}
\definecolor{bisque}{rgb}{1.0, 0.89, 0.77}
\definecolor{amber}{rgb}{1.0, 0.75, 0.0}
\definecolor{armygreen}{rgb}{0.29, 0.33, 0.13}
\definecolor{alizarin}{rgb}{0.82, 0.1, 0.26}
\definecolor{cadetblue}{rgb}{0.37, 0.62, 0.63}
\newcommand*{\TitleParbox}[1]{\parbox[c]{6cm}{\raggedright #1}}%
\newcommand{\python}{\texttt{python}}
\newcommand{\textoazul}[1]{\textcolor{blue}{#1}}
\newcommand{\azulfuerte}[1]{\textcolor{blue}{\textbf{#1}}}
\newcommand{\funcionazul}[1]{\textcolor{blue}{\textbf{\texttt{#1}}}}
%\normalfont
\usepackage{ccfonts}% http://ctan.org/pkg/{ccfonts}
\usepackage[T1]{fontenc}% http://ctan.or/pkg/fontenc
\renewcommand{\rmdefault}{cmr}% cmr = Computer Modern Roman
\usefonttheme[onlymath]{serif}
\linespread{1.3}
\newcounter{saveenumi}
\newcommand{\seti}{\setcounter{saveenumi}{\value{enumi}}}
\newcommand{\conti}{\setcounter{enumi}{\value{saveenumi}}}
\newcommand{\tikzmark}[1]{\tikz[remember picture] \node[coordinate] (#1) {#1};}

\usepackage{scalerel}[2016-12-29]
\def\stretchint#1{\vcenter{\hbox{\stretchto[440]{\displaystyle\int}{#1}}}}
\def\scaleint#1{\vcenter{\hbox{\scaleto[3ex]{\displaystyle\int}{#1}}}}
\def\bs{\mkern-12mu}

\newtheorem{teo}{}[section]
\usepackage{blkarray}

%reduce el tamaño de letra de la etiqueta equations
\makeatletter
\def\maketag@@@#1{\hbox{\m@th\normalfont\small#1}}
\makeatother

%se usa para la x en itemize
\newcommand{\xmark}{\text{\ding{55}}}

%\AtBeginDocument{\setlength{\tymin}{1em}}


\definecolor{myblue}{rgb}{.8, .8, 1}

\usepackage{empheq}

\newlength\mytemplen
\newsavebox\mytempbox

\makeatletter
\newcommand\mybluebox{%
    \@ifnextchar[%]
       {\@mybluebox}%
       {\@mybluebox[0pt]}}

\def\@mybluebox[#1]{%
    \@ifnextchar[%]
       {\@@mybluebox[#1]}%
       {\@@mybluebox[#1][0pt]}}

\def\@@mybluebox[#1][#2]#3{
    \sbox\mytempbox{#3}%
    \mytemplen\ht\mytempbox
    \advance\mytemplen #1\relax
    \ht\mytempbox\mytemplen
    \mytemplen\dp\mytempbox
    \advance\mytemplen #2\relax
    \dp\mytempbox\mytemplen
    \colorbox{myblue}{\hspace{1em}\usebox{\mytempbox}\hspace{1em}}}

\makeatother



%Se usa la plantilla Warsaw modificada con whale
\mode<presentation>
{
  \usetheme{Warsaw}
  \setbeamertemplate{headline}{}
  %\useoutertheme{infolines}
  \usecolortheme{whale}
  \setbeamercovered{invisible}
  

 \setbeamertemplate{section in toc}[sections numbered]
 \setbeamertemplate{subsection in toc}[subsections numbered]
 \setbeamertemplate{subsection in toc}{\leavevmode\leftskip=3.2em\rlap{\hskip-2em\inserttocsectionnumber.\inserttocsubsectionnumber}\inserttocsubsection\par}
% \setbeamercolor{section in toc}{fg=blue}
 \setbeamercolor{subsection in toc}{fg=blue}
 \setbeamerfont{subsection in toc}{size=\small}


\setbeamertemplate{navigation symbols}{}
\setbeamertemplate{caption}[numbered]
% \setbeamercolor{frametitle}{fg=yellow,bg=blue!70!white}
\setbeamercolor{section in head/foot}{bg=black, fg=white}
%\setbeamercolor{subsection in head/foot}{bg=gray!30,fg=black}
%\setbeamercolor{author in head/foot}{fg=yellow}
%\setbeamercolor{date in head/foot}{fg=blue}

%\mode<presentation>
%{
%  \usetheme{Warsaw}
%  \setbeamertemplate{headline}{}
%  %\useoutertheme{infolines}
%  \useoutertheme{default}
%  \setbeamercovered{invisible}
%  % or whatever (possibly just delete it)
%}
}


\usepackage{courier}
\usepackage{listingsutf8}
\usepackage{listings}
\usepackage{xcolor}
\usepackage{textcomp}
\usepackage{color}
\definecolor{deepblue}{rgb}{0,0,0.5}
\definecolor{brown}{rgb}{0.59, 0.29, 0.0}
\definecolor{OliveGreen}{rgb}{0,0.25,0}
% \usepackage{minted}

\DeclareCaptionFont{white}{\color{white}}
\DeclareCaptionFormat{listing}{\colorbox{gray}{\parbox{0.98\textwidth}{#1#2#3}}}
\captionsetup[lstlisting]{format=listing,labelfont=white,textfont=white}
\renewcommand{\lstlistingname}{Código}


\definecolor{Code}{rgb}{0,0,0}
\definecolor{Keywords}{rgb}{255,0,0}
\definecolor{Strings}{rgb}{255,0,255}
\definecolor{Comments}{rgb}{0,0,255}
\definecolor{Numbers}{rgb}{255,128,0}

\makeatletter

\newif\iffirstchar\firstchartrue
\newif\ifstartedbyadigit
\newif\ifprecededbyequalsign

\newcommand\processletter
{%
  \ifnum\lst@mode=\lst@Pmode%
    \iffirstchar%
        \global\startedbyadigitfalse%
      \fi
      \global\firstcharfalse%
    \fi
}

\newcommand\processdigit
{%
  \ifnum\lst@mode=\lst@Pmode%
      \iffirstchar%
        \global\startedbyadigittrue%
      \fi
      \global\firstcharfalse%
  \fi
}

\lst@AddToHook{OutputOther}%
{%
  \lst@IfLastOtherOneOf{=}
    {\global\precededbyequalsigntrue}
    {}%
}

\lst@AddToHook{Output}%
{%
  \ifprecededbyequalsign%
      \ifstartedbyadigit%
        \def\lst@thestyle{\color{orange}}%
      \fi
    \fi
  \global\firstchartrue%
  \global\startedbyadigitfalse%
  \global\precededbyequalsignfalse%
}

\lstset{ 
language=Python,                % choose the language of the code
basicstyle=\footnotesize\ttfamily,       % the size of the fonts that are used for the code
numbers=left,                   % where to put the line-numbers
numberstyle=\scriptsize,      % the size of the fonts that are used for the line-numbers
stepnumber=1,                   % the step between two line-numbers. If it is 1 each line will be numbered
numbersep=5pt,                  % how far the line-numbers are from the code
backgroundcolor=\color{white},  % choose the background color. You must add \usepackage{color}
showspaces=false,               % show spaces adding particular underscores
showstringspaces=false,         % underline spaces within strings
showtabs=false,                 % show tabs within strings adding particular underscores
frame=single,   		% adds a frame around the code
tabsize=2,  		% sets default tabsize to 2 spaces
captionpos=t,   		% sets the caption-position to bottom
breaklines=true,    	% sets automatic line breaking
breakatwhitespace=false,    % sets if automatic breaks should only happen at whitespace
escapeinside={\#},  % if you want to add a comment within your code
stringstyle =\color{OliveGreen},
%otherkeywords={{as}},             % Add keywords here
keywordstyle = \color{blue},
commentstyle = \color{black},
identifierstyle = \color{black},
literate=%
         {á}{{\'a}}1
         {é}{{\'e}}1
         {í}{{\'i}}1
         {ó}{{\'o}}1
         {ú}{{\'u}}1
%
%keywordstyle=\ttb\color{deepblue}
%fancyvrb = true,
}

\lstdefinestyle{FormattedNumber}{%
    literate={0}{{\textcolor{red}{0}}}{1}%
             {1}{{\textcolor{red}{1}}}{1}%
             {2}{{\textcolor{red}{2}}}{1}%
             {3}{{\textcolor{red}{3}}}{1}%
             {4}{{\textcolor{red}{4}}}{1}%
             {5}{{\textcolor{red}{5}}}{1}%
             {6}{{\textcolor{red}{6}}}{1}%
             {7}{{\textcolor{red}{7}}}{1}%
             {8}{{\textcolor{red}{8}}}{1}%
             {9}{{\textcolor{red}{9}}}{1}%
             {.0}{{\textcolor{red}{.0}}}{2}% Following is to ensure that only periods
             {.1}{{\textcolor{red}{.1}}}{2}% followed by a digit are changed.
             {.2}{{\textcolor{red}{.2}}}{2}%
             {.3}{{\textcolor{red}{.3}}}{2}%
             {.4}{{\textcolor{red}{.4}}}{2}%
             {.5}{{\textcolor{red}{.5}}}{2}%
             {.6}{{\textcolor{red}{.6}}}{2}%
             {.7}{{\textcolor{red}{.7}}}{2}%
             {.8}{{\textcolor{red}{.8}}}{2}%
             {.9}{{\textcolor{red}{.9}}}{2}%
             {\ }{{ }}{1}% handle the space
         ,%
          %mathescape=true
          escapeinside={__}
          }



\makeatletter

% \setbeamercolor{subsection in foot}{bg=blue!30!yellow, fg=red}
%\setbeamercolor{footlinecolor}{bg=black,fg=white}
\setbeamertemplate{footline}
{
  \leavevmode%
  \hbox{%
  \begin{beamercolorbox}[wd=.333333\paperwidth,ht=2.25ex,dp=1ex,center]{section in footline}%
    \usebeamerfont{section in foot} \insertsection
  \end{beamercolorbox}}%
  \begin{beamercolorbox}[wd=.333333\paperwidth,ht=2.25ex,dp=1ex,center]{subsection in foot}%
    \usebeamerfont{subsection in foot}  \insertsubsection
  \end{beamercolorbox}%
  \begin{beamercolorbox}[wd=.333333\paperwidth,ht=2.25ex,dp=1ex,right]{date in head/foot}%
    \usebeamerfont{date in head/foot} \insertshortdate{} \hspace*{2em}
    \insertframenumber{} / \inserttotalframenumber \hspace*{2ex} 
  \end{beamercolorbox}}%
  \vskip0pt%
\makeatother
\normalfont
\usepackage{ccfonts}% http://ctan.org/pkg/{ccfonts}
\usepackage[T1]{fontenc}% http://ctan.or/pkg/fontenc
\renewcommand{\rmdefault}{cmr}% cmr = Computer Modern Roman
\linespread{1.3}
\title{Solución al problema del satélite}
\subtitle{Curso de Física Computacional}
\author{M. en C. Gustavo Contreras Mayén}
\date{\today}
\institute{Facultad de Ciencias - UNAM}
\titlegraphic{\includegraphics[width=1.75cm]{Imagenes/escudo-facultad-ciencias.jpg}\hspace*{4.75cm}~%
   \includegraphics[width=1.75cm]{Imagenes/escudo-unam.jpg}
}
%\newcommand{\tikzmark}[1]{\tikz[overlay,remember picture] \node (#1) {};}
\newcommand{\DrawBox}[3][]{%
    \tikz[overlay,remember picture]{
    \draw[black,#1]
      ($(#2)+(-0.5em,2.0ex)$) rectangle
      ($(#3)+(0.75em,-0.75ex)$);}
}
\begin{document}
\maketitle
\fontsize{14}{14}\selectfont
\spanishdecimal{.}
\begin{frame}
\frametitle{El problema del satélite}
\begin{center}
\begin{tikzpicture}[font=\small]
	\draw (2,2) circle (2cm);
	\draw [dashed] (-0.4,2) -- (5,2);
	\draw [->] (2,2) -- node [midway, fill=white] {$R_{t}$} (1.8,4);
	\draw (5,2) circle (0.05);
	\draw (5.2,1.8) node {H}; 
	\draw [->] (5,2) -- node [near end, right] {$v_{0}$}(5,2.5);
	\draw [->] (2,2) -- node [midway, above] {$r$} (4,3.6);
	\draw (2.7,2) arc (0:45:5.9mm) node [right] {$\theta$};
	\draw [dashed] (5,2) arc (0:95:1.8);
\end{tikzpicture}
\end{center}
Un satélite se lanza desde una altitud $H=772$ km sobre el nivel del mar, con una velocidad inicial $v_{0}=6700$ $m/s$ en la dirección que se muestra.
\end{frame}
\begin{frame}
\frametitle{Sistema de EDO-2}
El conjunto de EDO que describen el movimiento del satélite son:
\[ \ddot{r} = r  \dot{\theta}^{2} - \dfrac{G M_{t}}{r^{2}}  \hspace{2cm} \ddot{\theta} = - \dfrac{2 \dot{r}\dot{\theta}}{r}\]
donde $r$ y $\theta$ son las coordenadas polares del satélite.
\end{frame}
\begin{frame}
\frametitle{Constantes del Problema}
Las constantes involucradas en las expresiones, son:
\begin{eqnarray*}
G &=& 6.672 \times 10^{-11} \mbox{ m}^{3} \mbox{kg}^{-1} \mbox{s}^{2} \\
M_{t} &=& 5.9742 \times 10^{24} \mbox{ kg, Masa de la Tierra} \\
R_{e} &=& 6378.14 \mbox{ km, radio de la Tierra al nivel del mar} 
\end{eqnarray*}
\end{frame}
\begin{frame}
\frametitle{Problemas a resolver}
\setbeamercolor{item projected}{bg=purple!80!white,fg=white}
\setbeamertemplate{enumerate items}[circle]
\begin{enumerate}[<+->]
\item Obtén el conjunto de EDO-1 y las condiciones iniciales del problema, de la forma $\dot{\mathbf{y}} = \mathbf{F}(x,\mathbf{y})$, $\mathbf{y}(0) = \mathbf{b}$.
\item Integra las EDO-1 en el tiempo en que se lanza el satélite y choca en su regreso a la Tierra.
\item Calcula el lugar del impacto, con $\theta$.
\end{enumerate}
\end{frame}
\begin{frame}
\frametitle{Inciso 1: Sistema de EDO-1}
Tenemos que
\[ \begin{split} G M_{t} &= (6.672 \times 10^{-11}) (5.9742 \times 10^{24})  = \\
 &= 3.9860 \times 10^{14} \mbox{m}^{3} \mbox{s}^{2} \end{split} \]
Ahora hacemos
\[ \mathbf{y} = \begin{bmatrix}
y_{0} \\
y_{1} \\
y_{2} \\
y_{3}
\end{bmatrix} =
\begin{bmatrix}
r \\
\dot{r} \\
\theta \\
\dot{\theta}
\end{bmatrix}
\]
\end{frame}
\begin{frame}
\frametitle{Inciso 1: Sistema de EDO-1}
Por lo que el conjunto equivalente de EDO-1, resulta ser
\[ \dot{\mathbf{y}} = \begin{bmatrix}
\dot{y}_{0} \\
\dot{y}_{1} \\
\dot{y}_{2} \\
\dot{y}_{3}
\end{bmatrix} =
\begin{bmatrix}
y_{1} \\
y_{0} y_{3}^{2} - 3.9860 \times 10^{14}/y_{0}^{2} \\
y_{3} \\
-2 y_{1} y_{3} / y_{0}
\end{bmatrix} \]
\end{frame}
\begin{frame}
\frametitle{Condiciones iniciales}
Las condiciones iniciales resultan
\fontsize{12}{12}\selectfont
\begin{eqnarray*}
r(0) &=& R_{t} + H = (6378.14+772) \times 10^{3} = 7.15014 \times 10^{6} \mbox{ m} \\
\dot{r}(0) &=& 0 \\
\theta(0) &=& 0 \\
\dot{\theta}(0) &=& \dfrac{v_{0}}{r(0)} = \dfrac{6700}{7.15014 \times 10^{6}} = 0.937045 \times 10^{-3} \mbox{ rad/s}
\end{eqnarray*}
\fontsize{14}{14}\selectfont
\pause
Por tanto
\[ \mathbf{y}(0) = 
\begin{bmatrix}
7.15014 \times 10^{6} \\
0 \\
0 \\
0.937045 \times 10^{3}
\end{bmatrix}
\]
\end{frame}
\begin{frame}
\frametitle{Inciso 2: Integración de las EDO-1}
Usaremos la función \funcionazul{integrate.odeint} para integrar el sistema de EDO-1.
\\
\bigskip
De lo visto en clase, se necesita la función \funcionazul{y, t} que incluye el sistema de EDO-1.
\end{frame}
\begin{frame}
\frametitle{Inciso 2: Integración de las EDO-1}
El período de integración, es decir, el valor del tiempo que tarda en caer el satélite y choca con la Tierra, hay que estimarlo tentativamente.
\end{frame}
\begin{frame}[plain, allowframebreaks, fragile]
\begin{lstlisting}[caption=Solución con python, style=FormattedNumber, basicstyle=\linespread{1.1}\ttfamily=\small, columns=fullflexible]
def F(y, t):
    F = zeros((4), dtype='float64')
    F[_0_] = y[_1_]
    F[_1_] = y[_0_] * (y[_3_]**2) - 3.9860e14 / (y[_0_]**2)
    F[_2_] = y[_3_]
    F[_3_] = -2.0 * y[_1_] * y[_3_]/y[_0_]
    return F

t = np.linspace(0, 1200, 100)
y_0_ = np.array([7.1514e6, 0., 0., 0.937045e-3])
y_1_ = odeint(F, y_0_, t)

print('tiempo \t altura \t velocidad')
print('-'*30)

for i in range(len(t)):
    print('{0:4.4f} \t {1:1.6e} \t {2:1.6e}'.format(t[i], y_1_[i][_0_], y_1_[i][_1_]))
\end{lstlisting}
\end{frame}
\begin{frame}[fragile]
\frametitle{Resultados}
\fontsize{10}{10}\selectfont
\begin{center}
\hspace*{-1.25cm}
\begin{tabular}{c | c | c | c | c }
tiempo & altura ($r$) & velocidad ($\dot{r}$) & $\theta$ & $\dot{\theta}$ \\ \hline
$0.0000$ & $7.151400e+06$ & $0.000000e+00$ & $0.000000e+00$ & $9.370450e-04$\\ \hline
$12.1212$ & $7.151289e+06$ & $-1.835853e+01$ & $1.135824e-02$ & $9.370742e-04$ \\ \hline
$24.2424$ & $7.150955e+06$ & $-3.671583e+01$ & $2.271718e-02$ & $9.371616e-04$ \\ \hline
\ldots & \ldots & \ldots & \ldots & \ldots \\ \hline
$1030.3030$ & $6.386228e+06$ & $-1.396748e+03$ & $1.043431e+00$ & $1.175043e-03$ \\ \hline
$1042.4242$ & $6.369227e+06$ & $-1.408222e+03$ & $1.057712e+00$ & $1.181324e-03$ \\ \hline
\ldots & \ldots & \ldots & \ldots & \ldots \\ \hline
\end{tabular}
\end{center}
\end{frame}
\begin{frame}[fragile]
\frametitle{Intervalo de tiempo al choque}
El satélite choca con la Tierra cuando $r$ es igual a $R_{t}=6.37814 \times 10^{6}$ m. 
\\
\bigskip
De los resultados, vemos que esto ocurre entre el tiempo $t= 1030$ y $1042$ segundos.
\end{frame}
\begin{frame}[fragile]
\frametitle{Intervalo de tiempo al choque}
\fontsize{10}{10}\selectfont
\begin{center}
\hspace*{-1.25cm}
\begin{tabular}{c | c | c | c | c}
tiempo & altura ($r$) & velocidad ($\dot{r}$) & $\theta$ & $\dot{\theta}$ \\ \hline
$0.0000$ & $7.151400e+06$ & $0.000000e+00$ & $0.000000e+00$ & $9.370450e-04$\\ \hline
$12.1212$ & $7.151289e+06$ & $-1.835853e+01$ & $1.135824e-02$ & $9.370742e-04$ \\ \hline
$24.2424$ & $7.150955e+06$ & $-3.671583e+01$ & $2.271718e-02$ & $9.371616e-04$ \\ \hline
\ldots & \ldots & \ldots & \ldots & \ldots \\ \hline
\tikzmark{top left 1}$1030.3030$ & $6.386228e+06$ & $-1.396748e+03$ & $1.043431e+00$ & $1.175043e-03$ \\ \hline
$1042.4242$ & $6.369227e+06$ & $-1.408222e+03$ & $1.057712e+00$ & $1.181324e-03$ \tikzmark{bottom right 1} \\ \hline
\ldots & \ldots & \ldots & \ldots & \ldots \\ \hline
\end{tabular}
\DrawBox[ultra thick, red]{top left 1}{bottom right 1}
\end{center}
\end{frame}
\begin{frame}
\frametitle{Intervalo de tiempo al choque}
Un valor de $t$ más preciso, lo podemos obtener mediante una interpolación  polinomial, pero si no queremos una precisión alta, con una interpolación lineal, bastará.
\\
\medskip
\pause
Consideremos que $1030.3030 + \Delta t$ el tiempo para el impacto, por lo que escribimos
\[ r(1030.3030 + \Delta t) = R_{t} \]
\end{frame}
\begin{frame}[fragile]
\frametitle{Intervalo de tiempo al choque}
Desarrollando $r$ con dos términos de la serie de Taylor, tenemos que
\[ \begin{split} 
r(1030.3030) + \dot{r}(1030.3030) \Delta t &= R_{t} \\
6.386228 \times 10^{6} + (-1.385053 \times 10^{3}) \Delta t &= R_{t}
\end{split} \]
que al despejar, 
\[ \Delta t =  \dfrac{R_{t} - 6.386228 \times 10^{6}}{-1.385053 \times 10^{3}} =  \SI{5.8394}{\second} \]
\pause
\textcolor{blue}{Por lo que el tiempo de impacto es $1036.1424$ segundos}.
\end{frame}
\begin{frame}
\frametitle{Inciso 3) Coordenada de impacto}
La coordenada $\theta$ del impacto, la podemos calcular de una manera similar, para ello desarrollamos dos términos de la serie de Taylor:
\[ \begin{split}
\theta(1030.3030 + \Delta t) &= \theta(1030.3030) + \dot{\theta}(1030.3030)\Delta t \\
 &= 1.043431 + (1.175043 \times 10^{-3} )(5.8394) \\
 &= 1.5029 \mbox{ rad} = 60.18^{\circ}
\end{split} \]
\end{frame}
\end{document}
